\chapter[Trade Carries Objects and Ideas]{Trade Lets Objects Travel, and Ideas Too}
\section{A Knife That Should Not Exist}
Visit a museum's bronze gallery: a small knife or arrowhead lies under glass, green corrosion like dried moss. Its label says around four thousand years old.

No trains, steamships, roads, or widely used chariots then. Look one step further: this object should not exist.

Bronze is not excavated as bronze. Copper and tin ores emerge from earth, but bronze is a recipe: roughly nine parts copper to one tin, melted and cooled into something harder than copper and easier to shape than stone. Nature seems perversely to separate sources. Copper-bearing places often lack tin; tin-bearing places lack copper. Copper is fairly widespread in mountains, tin much rarer, with few major ancient sources, mostly remote. Near Eastern bronze workers' tin origins remain disputed, even earning a name, the tin problem.

Before becoming a knife, this object was therefore a journey. Copper left one mountain, tin another direction perhaps hundreds or thousands of kilometres away, meeting at a furnace before the artisan lit a fire. Motionless under glass, its body freezes at least two routes, many hands, repeated bargaining. Made of metal, it is also made of roads.

Look further. Ur's royal tombs in today's Iraq yielded lapis beads over four millennia old, deep blue as midnight, whose nearest major mines lay over two thousand kilometres away in northeastern Afghanistan's Badakhshan. Pacific obsidian blades chemically trace to volcanic islands thousands of kilometres distant. Tang tombs near Xi'an hold Persian silver; East African beaches hold Song Chinese porcelain.

Objects cannot walk. How did they travel, who carried them, and what else did carriers take or leave? Trade carries objects and ideas, and travel weighs more than it sounds.

\section{A Tablet in Ur}
Enter a scene.

Time: over thirty-seven centuries ago, Old Babylonian times. Place: lower-Euphrates Ur, smelling of river water, drying fish, burning camel dung. Its busy harbour warehouses pile reeds, oil jars, date sacks against mudbrick while scales rise and fall.

Merchant Nanni complains in a courtyard. He paid for high-grade copper ingots delivered by water, probably on the major Magan--Dilmun--Ur route: metal from eastern mountains around today's Oman, Magan, shipped through Dilmun, roughly Bahrain, a Gulf hub of transshipment, weighing, taxes, then upriver. Inspection reveals poor quality, worsened by insult to his collecting servant.

He responds as Mesopotamians do to serious affairs: summon a scribe and dictate a stern complaint. An angled reed presses cuneiform into wet clay; the dried letter goes to copper trader Ea-nasir.

Millennia later archaeologists excavated it; it now rests in the British Museum. The world's oldest bad review, people joke. More than amusement, it preserves a scene containing quality standards, long waterways and intermediate ports, agents and servants, credit and disputes, and written demands for redress.

Stay in that courtyard, smell the river, inspect the infuriating ingots. Everything following answers questions hidden there.

\section{Is Trade Merely Moving Things?}
Set out the conflict before answering.

A plain account, almost tautological, calls trade transport. Abundant cheap goods in A move to scarce expensive B; merchants take the difference. Nanni's story becomes copper reaching a city, prices multiplying, profits taken, metal supplied, everyone home.

A few further questions spring leaks.

First: why pay that difference? Lapis may be merely attractive in Badakhshan, but at Ur temples and royals pay vastly more for divine gifts and insignia. Distance, not hardness or utility, creates value. Difficult access displays owners' capacity. Trade sells remoteness too.

Second: who needs whom more? Ur grows food but lacks copper and tin, depending on outside metal for weapons, tools, sacred vessels. Mountain miners and herders could live without its grain, cloth, orders, seemingly making the city need trade more. Yet miners accustomed to exchange cannot simply resume old lives. Dependence grows both ways and becomes hard to dismantle. Who holds whose vulnerable point?

Third: do objects change en route? Tin begins as ore, becomes alloy ingredient, weapon, battlefield power. Lapis becomes a god's eyes in a temple. Each hand and place rewrites meaning. Does trade move objects or meanings?

Fourth, quietest: what travels alongside them? Nanni's cuneiform arose from bookkeeping and receipts, earliest writing counting copper rather than poems. Caravans carry languages, alphabets, stories, deities, pathogens attached to goods. Invisible passengers buy no tickets, undergo no weighing, pay no tax, yet change worlds as much as metals.

None receives an immediate courtyard answer. Choose provisionally: transport, profit, or something else as trade's essence? Are endpoints equal partners, or must one eventually grip the other?

\section[Kings, Traders, Interpreters, Miners]{Kings, Merchants, Interpreters, and People in the Mines}
Many people occupy a route, seeing different trades. Hear several.

\textbf{Kings and temples: routes are lifelines.} Imported copper, tin, timber, stone are strategic. Interrupted routes deprive armies of weapons, temples of adornment, kingship of display. Near Eastern rulers protect merchants and sign treaties while trading themselves: court and temple are major buyers, stores, lenders. This is national security, not free markets. Route control grips others' throats. Millennia later, oil and chips replace copper and tin within unchanged logic.

\textbf{Merchants: profit from prices and trust.} Cheapness, prices, safe roads, recoverable payments matter. Yet trust, not goods, is long-distance trade's scarcest resource. Entrust valuable copper to a ship and agent for months, without courts litigating internationally: why expect delivery? People and relational networks answer. We will examine that true engine later.

\textbf{Interpreters and diasporas: living bridges.} Goods cannot speak across language boundaries. More important still, settled compatriot communities understand both home rules and local markets. Sogdians from around Samarkand in today's Uzbekistan famously connected their homeland to Chang'an and Luoyang; their language became a Silk Road lingua franca. Unsent family letters survived in a Han watchtower near Dunhuang. Song--Yuan Guangzhou and Quanzhou had foreign quarters where Arab and Persian traders lived, traded, worshipped, sometimes for generations, leaving bilingual Arabic-Chinese tombstones. Not merely carriers, these people embody the network.

\textbf{Mine workers and caravan labourers: no chance to speak.} Behind travelling objects stand unrecorded miners, camel drivers, dock porters. Lapis's first kilometre out of Badakhshan rode a human back. Saharan salt exchanged for gold, but cutters worked in desiccating heat. Glossy object biographies omit them, though without them page one cannot begin. Including them in trade's glory is honesty, not spoiling the mood.

\textbf{Customers: Nanni's bad reviews.} Do not forget the first complainer. Ordinary people care less where routes reach than whether copper is good and cloth worth its price. Grand trade histories finally land at innumerable inspections like his.

Another position deserves respect despite frequent ridicule: \textbf{self-sufficiency}, avoiding outside exchange where possible.

Give it its strongest case because today's trade-is-always-good default treats opponents unfairly. Four arguments follow. Dependence exposes food and weapon materials to cutoffs, embargoes, blockades. Routes carry pathogens alongside goods, with bloody evidence ahead. Cheap superior imports destroy local crafts; after generations, skills cannot be recovered. Profits widen inequality between those near routes and those distant. Ming--Qing maritime bans and Tokugawa Japan's two-century closure partly followed the reasoning that smaller openings mean safety. Reject their conclusion if you wish, but all four concerns are real and recur in new forms.

\section[Thinking Bridge: An Object's Biography]{A Thinking Bridge: Write an Object's Biography}
Trade changing the world needs a handheld tool: \textbf{object biography}. As with a person, trace its life through five questions.

\textbf{First: Where was it born?} Begin with raw-material origins, not object classification. Bronze has at least two birthplaces, copper and tin mines; cola ingredients may come from five or six countries. Origins reveal an invisible map.

\textbf{Second: Which route?} Sea or land, intermediate ports, taxes, languages? Choosing this route rather than another often encodes geography, politics, technology's whole secret.

\textbf{Third: Whose hands received it?} One person rarely carries long-distance goods end to end. Ancient Silk Roads were relays, not express couriers, changing merchants at successive markets. This deserves emphasis:

A \textbf{particularly easy-to-misread counterexample}: distant-sourced Pacific obsidian once encouraged explanations of migration or conquest. More ordinary relay exchange may suffice. Stone travelling three thousand kilometres does not require anyone travelling that far. Dozens of handovers can move it in small stages while owners remain lifelong in one bay. Object distances often exceed human distances. Inferring distant arrivals of people from imported things is a common archaeological-news mistake.

\textbf{Fourth: What did it become at destination?} Most exciting, because travel changes an object's heart. Lapis shifts from attractive stone to divine material. Chinese clothing silk becomes Roman luxury debated for prohibition; thin male garments once provoked outraged senators. Local needs reinterpret every arrival. Trade repeatedly renames rather than merely relocates.

\textbf{Fifth: Where did it die, and what record remains?} Royal grave, seabed, construction rubble, ledger, inscription? Endings reveal what an age valued or discarded.

Practise with Saharan salt. Birth: deep-desert deposits such as Taghaza, cut into regular slabs. Route: camel caravans cross for weeks through supplying oases. Hands: guides, oasis merchants, West African buyers. Transformation: where salt is rarer than gold, slabs buy gold dust, not legendary equal-weight exchange but genuinely hard currency, every soup grain money carried hundreds of kilometres. Death: dissolved in countless meals, shadows remaining in accounts and oral traditions. Five questions reveal a trans-Saharan network. You learn eyes for cross-regional history, not merely salt's story.

\subsection[Invisible Passengers]{Invisible Passengers: Alphabets, Stories, and Pathogens}
Push the tool further: unboxed cargo hitchhikes farther than any commodity.

\textbf{Travelling letters.} This book's Chinese original uses locally developed characters, but English's twenty-six letters trace ancestry to the eastern Mediterranean over three millennia ago. Phoenician seafaring merchants used a twenty-odd-sign alphabet more convenient for daily accounts than hundreds of pictographs. Ships carried it west; Greeks added vowels to consonant signs; Romans received it as the Latin alphabet, eventually spreading across much of earth. Keyboard letters descend from Mediterranean merchants' shorthand. Scripts travel through costs as well as armies and dynasties: easy learning and writing save traders money, and savings find their own way.

\textbf{Travelling stories.} Tang writer Duan Chengshi's Miscellaneous Morsels from Youyang records Ye Xian, abused by her stepmother, whose magical fish is killed but whose powerful bones grant golden shoes. A local ruler finds one, searches for its owner, marries her. Cinderella readers recognise striking resemblance centuries before European written versions. Did it travel? Honestly, unknown. Abused orphan, magical helper, identifying token can arise independently like paddles. Yet caravans and migrants truly carry stories: Indian Jatakas followed Buddhist traders into Central Asia and China; Arab sailors' tales left mixed descendants on East African coasts. Independent invention versus diffusion remains debated. Unusual shared details matter: common skeletons may coincide, matching obscure particulars suggest kinship.

\textbf{Travelling pathogens.} The heaviest passengers. Fourteenth-century plague followed Eurasian routes from Central Asian caravans to Black Sea Kaffa, then Genoese ships full of goods and rats into the Mediterranean, sweeping Europe within years. Roughly one-third to half of Europeans died, a broadly accepted order of magnitude. Rats and fleas bought no tickets while spices were weighed and taxed. Efficient pipes do not select their contents. Later European ships brought smallpox to immunologically unprepared Indigenous Americans; merchant vessels and pathogens followed conquerors' swords. Celebrating connection without these costs falsifies history.

\section[Zhang Qian Finds Goods from Home]{Zhang Qian Finds Goods from Home in Bactria}
Read primary material from "Accounts of Dayuan" in Records of the Grand Historian. Han envoy Zhang Qian, held by the Xiongnu for over a decade during western missions, eventually opened routes through the Hexi Corridor to Central Asia. Returning to Chang'an, he reported to Emperor Wu:

\begin{quote}
In Daxia I saw Qiong bamboo staffs and Shu cloth. I asked their origin. The people replied: "Our merchants go to Shendu to buy them."
\end{quote}
Meaning: in Bactria, around today's Afghanistan, he saw staffs from Qiong and cloth from Sichuan's Shu. Local merchants had purchased them in India, Shendu.

Pause at his surprise. The central dynasty's far-west traveller, already immensely distant from home, finds Sichuan goods in an Afghan market. They bypassed his newly official northern route, travelling south through today's Yunnan mountains into India, then successive traders to Bactria. The Han court did not know this route existed.

That surprise is the document's value: \textbf{merchants' feet traced roads before official maps did}. Zhang's opening of the west fills textbooks, yet Shu cloth reveals states often stamp, garrison, and toll pre-existing popular networks. Antlike civilian exchange precedes banner-filled embassies. Do not reverse the sequence.

At the network's western end, over a century after Zhang, Pliny the Elder's Natural History complains of immense annual Roman silver outflows for eastern silk, spices, pearls, estimating tens of millions of sesterces. From moral high ground he laments men wearing silk and hard wealth entering others' pockets.

Together, an eastern envoy marvels at home goods far away while a western scholar mourns home money departing. Neither sees the whole route; nobody could walk it all. Each encountered fragment truly surprises or hurts. Cross-regional history normally offers local views everywhere shadowed by distance.

\section[Why Trade Routes Survive Millennia]{Why Can a Trade Route Keep Working for Two Thousand Years?}
We now have enough components to compare explanations. Why did the broad Eurasian Silk Road network persist intermittently for millennia through dynastic turnover, imperial collapse, epidemic? Separate evidence of enduring movement, explanations of persistence, participants' views, Pliny's moral decay versus Bactrian business, and our evaluation of connection. We compare explanations.

\textbf{Explanation 1: Price differences attract: economics.}

Water runs downhill, goods toward high prices. Silk affordable in China commands astonishing Roman prices; ordinary local spices become silver-priced Mediterranean luxuries. Gaps covering transport, loss, tariffs, bribes, mortal risk continually generate merchants. Bans and blockades merely add costs. Strong and simple, this explains routes regrowing after war: price gravity remains. But merchants are not water. Why Sogdians, Armenians, Arabs rather than similarly placed others? Turning ubiquitous gaps into money requires special equipment, relationships rather than camels.

\textbf{Explanation 2: Empires build roads: politics.}

Trade rests on order, not spontaneous nature. Distance matters less than bandits, barriers, incompatible money, unenforced contracts. Empires supply order as a public good: Han protectorates, watchtowers, relay stations; Persia's Royal Road; Mongol routes from Black Sea to Pacific enabling thirteenth- and fourteenth-century Europeans to reach China. Fair scales and effective passes belong to imperial peace. This explains prosperity's coincidence with stability. But trade survives collapse through detours, sea routes, fragmentation. Empires also charge and constrain, alternately protecting merchants and fleecing them. Crediting them with all trade resembles crediting dykes with rivers' existence.

\textbf{Explanation 3: Trust's scaffolding: social networks.}

Return to price theory's missing equipment. Nanni commits money to copper, vessels, agents, waits months without international courts or bank transfers, with little beyond stern letters against default. Trade works through partners unable to escape reputational consequences: relatives, compatriots, fellow worshippers. Samarkand-centred Sogdian kin and trading stations stretch to Chang'an. Chinese default ruins a whole family's reputation along the route, their only immovable asset, firmer than contracts. Mediterranean Jewish, Indian Ocean Arab and Indian, later Southeast Asian Fujianese and Cantonese networks share the mechanism: \textbf{diasporas are ancient letters of credit}. This explains which groups succeed and how trade survives imperial gaps. But inward scaffolds become outward barriers, excluding outsiders and rival networks, vulnerable to monopolisation and internal strife. Trust builds roads and boundaries.

Each holds truth. Prices supply gravity, empires surface, trust wheels; lacking any, vehicles struggle. Not mere compromise, but warning: theories explaining two millennia alone have probably discarded evidence.

\section[Further Challenge: Connection Locks In]{A Deeper Challenge: Connection's Other Face Is Lock-In}
Our strongest theory asks not why routes exist, but \textbf{why connection enriches some regions while impoverishing others, instead of benefiting everyone together}.

From the 1970s, Immanuel Wallerstein's multivolume The Modern World-System argues: \textbf{analyse the whole trading network rather than separate national boxes; the network generates its own division of labour}.

\textbf{Core} regions control capital, technology, prices, profitable processing, finance, rules. \textbf{Peripheries} supply raw materials and cheap labour, ore, spices, timber, effort, earning thin margins at prices others set. \textbf{Semiperipheries} serve cores while controlling outer regions, buffering between them. These positions reflect neither effort nor intelligence. Once assigned, the network can \textbf{lock} regions in: raw-material profits rationally encourage more extraction rather than competing industrialisation. Ports, roads, institutions orient around export, making reversal harder. Connection's other face is dependence.

Caribbean sugar islands illustrate it: land could grow other things, but colonial systems made sugar workshops, organising soil, slavery, ports, shipping around cane until falling prices revealed few alternatives. Spice Islands' cloves and nutmeg multiplied value in Europe, with ship and market controllers taking most, cultivators little. Individual exchanges appear voluntary and equivalent, yet \textbf{equivalent products stand on unequal positions}. Fair exchange can coexist with unfair division of labour, the theory's sharpest cut.

Weigh substantial criticisms. First, fatalism: East Asian economies moved from raw materials to advanced manufacturing within generations. Lock-in is not life imprisonment. Second, passive peripheries: geographically peripheral Sogdians actively embodied connection as hubs and intermediaries rather than simply being connected. Third, broad positions overlook miners, merchants, weavers, enslaved people living radically differently inside one network. Theory is map, not sentence. It describes terrain; choosing paths and evaluating them remains yours.

\section{A World Journey in Your Pocket}
The ancient problem sits in your pocket.

Write your phone's short biography. Much battery cobalt comes from Congolese mines whose conditions receive sustained scrutiny; copper and lithium may come from Chile and Australia; chips probably from Taiwan using Dutch lithography equipment assembled from many countries' parts. Components fly to Chinese or Vietnamese assembly, then containers sail to your city. Its prior travel may exceed your lifetime's. Like ancient Ur, the land beneath you cannot alone make what you hold.

The network stays invisible until knotted. In March 2021, Ever Given grounded sideways in Suez, blocking for six days a waterway carrying over a tenth of global maritime trade. Hundreds of ships queued; factories worldwide counted inventory. Six days and a narrow channel made supply chains bodily experience. Earlier pandemic shortages of masks, ventilator parts, then automotive chips showed concentrating efficiency also concentrates risk. Self-sufficiency's four arguments return as supply-chain security, de-risking, reshoring: vulnerable dependence, indiscriminate pipes, lost crafts, widening inequalities.

Ideas travel faster too. Buddhist motifs once took centuries by caravan; memes, recipes, dances circle earth in days. Classmates listen across continents and watch dramas across hemispheres. Milk-tea ingredients, trainer factories, game servers each add a new segment to old routes.

\section[Your Turn: Connection or Dependence?]{Your Turn: Where Does Connection Become Dependence?}
Return to the bronze knife.

Its maker four millennia ago did not know he participated in something spanning two or three centuries and thousands of kilometres. Congolese miners, Taiwanese wafer engineers, port crane drivers today rarely think of themselves as world-system positions. Yet network, connection, lock-in remain.

Our question: \textbf{where does connection become dependence?}

Answer with reasons after weighing both sides.

For connection: specialisation makes goods richer and cheaper; routes carry alphabets, stories, technologies; without them, no phone, textbook, perhaps breakfast sugar. Everyone has felt broken chains since 2020.

For caution: dependence hands others vulnerabilities; pipes carry pathogens and domination; equal exchange may conceal unequal positions, with extraction roles trapping generations. Complaining customers have persisted since Nanni's millennia-old letter.

Closure also costs: maritime bans and Tokugawa restrictions bought security with stagnation and missed chances. Complete openness can destroy crafts and communities with cheap imports. Advance and retreat both bring bills; recipients differ.

Entrusting food, energy, chips to global networks: rational specialisation or exposed vulnerability? If a cheap possession relies on distant people locked into mining, would you pay more? How much suffices?

No standard answer. Next news of supply chains, trade wars, cutoffs will reveal not abstract games but a four-thousand-year road carrying copper, tin, salt, spices, alphabets, stories, pathogens, and the glowing screen in your hand.
